% Laget automatisk av elevark/lag_elevark.py fra sim/luftmotstand.html. Endringer her blir overskrevet.
\documentclass[11pt]{article}
\usepackage{elevark}
\begin{document}
\elevarkhode{Fysikk 1 · Krefter}{Fall med luftmotstand}{https://fysikklab.no/sim/luftmotstand.html}{Simuleringen}

En muffinsform er lett og har stor flate, så luftmotstanden blir fort like stor som tyngden. Da faller den med konstant fart, terminalfarten. Hva skjer når du stabler flere former i hverandre? Formen blir den samme, men massen øker.

\section*{Oppgaver}

\begin{oppgaveliste}
\oppgave \fase{gjett}{Gjett}Slipp én form. Hvordan ser fartsgrafen ut? Hva skjer med luftmotstanden mens farten øker?
\svar
\oppgave \fase{forklar}{Forklar}Forklar med Newtons lover hvorfor formen etter hvert faller med konstant fart. Hva er summen av kreftene da?
\svar
\oppgave \fase{test}{Regn}Sett $L = mg$ når formen har terminalfart, og vis at $v_{\mathrm{t}} = \sqrt{mg/k}$. Regn ut terminalfarten for én form, og sjekk.
\svar
\oppgave \fase{gjett}{Gjett}Hva skjer med terminalfarten hvis du stabler fire former? Blir den fire ganger så stor? Test.
\svar
\oppgave \fase{test}{Undersøk}Slipp 1, 2, 3, 4 og 5 former, og les av terminalfarten hver gang. Lag en graf av $v_{\mathrm{t}}^{2}$ mot antall former. Hva viser grafen?
\svar
\oppgave \fase{test}{Undersøk}Slå på «Uten luftmotstand». Hvor mye kortere tid tar fallet? Hvor fort ville formen truffet gulvet?
\svar
\oppgave \fase{utfordring}{Utfordring}En fallskjermhopper har terminalfart omtrent 55~m/s før skjermen utløses, og omtrent 5~m/s etterpå. Hvor mange ganger større blir $k$ når skjermen utløses?
\svar
\end{oppgaveliste}

\end{document}
